% Weighted refinement; path_stability independently audited all steps and constants.
\section{A closure-mass version of the cut estimate}
\label{sec:closure-mass-cuts}

The degree bound and the local coordinate rank play different roles.
Keeping them separate gives a weighted bound with polynomial dependence on the maximum boundary size.

\begin{lemma}[Degree and local rank separately]\label{lem:degree-rank-cut-estimate}
Let $Q$ be connected and loopless, with $k$ distinct marked vertices.
Suppose their degrees are at most $d\ge2$ and their local coordinate spaces have dimension at most $s\ge1$.
If $Z$ counts the closed marked vertices in a uniform even subgraph, then
\[
 \mathbb E2^{-Z}\le\frac12+2\exp\!\left(-\frac{k}{8d2^{s+1}}\right).
\]
\end{lemma}

\begin{proof}
A zero local space gives an expectation at most one half.
Otherwise put $T=2^{s+1}$.
If $T$ marked local spaces form a direct sum, their closure indicators are independent.
Each has closure probability at least $2^{-s}$, so the expectation is at most
$(1-2^{-s-1})^T\le e^{-1}<1/2$.

Otherwise choose a maximal direct-sum family $S$ with $t<T$ vertices.
Every other marked local space meets $U=\sum_{v\in S}L_v$ nontrivially.
For $k\ge2(5d+1)T$, the set $S$ is nonempty and proper, and its boundary has size $B\le dt$.
Lemma~\ref{lem:weighted-block-conditioning}, with $\eta_v=1/2$, gives
\[
 \mathbb E2^{-Z}-\frac12
 \le\frac12\exp\!\left(-2^{-s-1}(k-t-5B)\right)
 \le\frac12\exp\!\left((5d+1)-\frac{k}{T}\right)
 \le\frac12 e^{-k/(2T)}.
\]
This is stronger than the stated bound.
For $k<2(5d+1)T$, one has $k/(8dT)<11/8<\ln4$.
The stated right side then exceeds one.
\end{proof}

\begin{corollary}[Total closure mass]\label{cor:closure-mass-cut-estimate}
Let $Q$ be connected and loopless, with distinct marked vertices of degree at most $d\ge2$.
For a uniform even subgraph, put
\[
 p_v=\Pr(v\text{ is closed}),\qquad \mu=\sum_{v\text{ marked}}p_v.
\]
Then
\begin{equation}\label{eq:closure-mass-cut-estimate}
 \mathbb E2^{-Z}
 \le\frac12+2\exp\!\left(-\frac{\mu}{16d(d-1)}\right).
\end{equation}
The same estimate bounds simultaneous forest restrictions in disjoint connected cyclic regions,
where $p_v$ is the probability that the corresponding region boundary is closed.
\end{corollary}

\begin{proof}
If some local space has dimension zero, the expectation is already at most one half.
Otherwise each dimension belongs to $\{1,\ldots,d-1\}$.
For each $s$, let $k_s$ count marked vertices of local dimension exactly $s$.
Their closure mass is $\mu_s=k_s2^{-s}$.
Some $s$ satisfies $\mu_s\ge\mu/(d-1)$.
Ignore the other marked factors and apply Lemma~\ref{lem:degree-rank-cut-estimate} to this class.
Its exponent is
\[
 \frac{k_s}{8d2^{s+1}}=\frac{\mu_s}{16d}
 \ge\frac{\mu}{16d(d-1)}.
\]
This proves~\eqref{eq:closure-mass-cut-estimate}.
The cyclic-region statement follows from the quotient and internal-cycle-space reduction.
\end{proof}

This bound retains the actual closure mass, which can be much larger than the estimate obtained from the maximum degree alone.
The regions must still be vertex-disjoint.
The estimate does not yet control an overlapping family of moderate-length cycles.
